\documentclass[10pt,a4paper]{amsart}
\usepackage{amsmath,amssymb,amsthm,graphicx}
% Retained minimal preamble for the Stacks Project chapter more-morphisms.tex (More on Morphisms).
% Verbatim excerpt of the chapter preamble: Xy-pic setup (native preamble lines 14-18) and the native
% theorem declarations with the mathematical macros (native preamble lines 154-217). Copied mechanically;
% nothing added, removed or redefined. The wrapper class and the amsmath/amssymb/amsthm/graphicx base
% packages are supplied by the assembler, so the class/externaldocument part of the native preamble is omitted.
% For commutative diagrams we use Xy-pic
\usepackage[all]{xy}

% We use 2cell for 2-commutative diagrams.
\xyoption{2cell}
% Theorem environments.
%
\theoremstyle{plain}
\newtheorem{theorem}[subsection]{Theorem}
\newtheorem{proposition}[subsection]{Proposition}
\newtheorem{lemma}[subsection]{Lemma}

\theoremstyle{definition}
\newtheorem{definition}[subsection]{Definition}
\newtheorem{example}[subsection]{Example}
\newtheorem{exercise}[subsection]{Exercise}
\newtheorem{situation}[subsection]{Situation}

\theoremstyle{remark}
\newtheorem{remark}[subsection]{Remark}
\newtheorem{remarks}[subsection]{Remarks}

\numberwithin{equation}{subsection}

% Macros
%
\def\lim{\mathop{\mathrm{lim}}\nolimits}
\def\colim{\mathop{\mathrm{colim}}\nolimits}
\def\Spec{\mathop{\mathrm{Spec}}}
\def\Hom{\mathop{\mathrm{Hom}}\nolimits}
\def\Ext{\mathop{\mathrm{Ext}}\nolimits}
\def\SheafHom{\mathop{\mathcal{H}\!\mathit{om}}\nolimits}
\def\SheafExt{\mathop{\mathcal{E}\!\mathit{xt}}\nolimits}
\def\Sch{\mathit{Sch}}
\def\Mor{\mathop{\mathrm{Mor}}\nolimits}
\def\Ob{\mathop{\mathrm{Ob}}\nolimits}
\def\Sh{\mathop{\mathit{Sh}}\nolimits}
\def\NL{\mathop{N\!L}\nolimits}
\def\CH{\mathop{\mathrm{CH}}\nolimits}
\def\proetale{{pro\text{-}\acute{e}tale}}
\def\etale{{\acute{e}tale}}
\def\QCoh{\mathit{QCoh}}
\def\Ker{\mathop{\mathrm{Ker}}}
\def\Im{\mathop{\mathrm{Im}}}
\def\Coker{\mathop{\mathrm{Coker}}}
\def\Coim{\mathop{\mathrm{Coim}}}

% Boxtimes
%
\DeclareMathSymbol{\boxtimes}{\mathbin}{AMSa}{"02}

%
% Macros for moduli stacks/spaces
%
\def\QCohstack{\mathcal{QC}\!\mathit{oh}}
\def\Cohstack{\mathcal{C}\!\mathit{oh}}
\def\Spacesstack{\mathcal{S}\!\mathit{paces}}
\def\Quotfunctor{\mathrm{Quot}}
\def\Hilbfunctor{\mathrm{Hilb}}
\def\Curvesstack{\mathcal{C}\!\mathit{urves}}
\def\Polarizedstack{\mathcal{P}\!\mathit{olarized}}
\def\Complexesstack{\mathcal{C}\!\mathit{omplexes}}
% \Pic is the operator that assigns to X its picard group, usage \Pic(X)
% \Picardstack_{X/B} denotes the Picard stack of X over B
% \Picardfunctor_{X/B} denotes the Picard functor of X over B
\def\Pic{\mathop{\mathrm{Pic}}\nolimits}
\def\Picardstack{\mathcal{P}\!\mathit{ic}}
\def\Picardfunctor{\mathrm{Pic}}
\def\Deformationcategory{\mathcal{D}\!\mathit{ef}}
\usepackage{cleveref}
\crefname{originalreference}{original reference}{original references}
\Crefname{originalreference}{Original reference}{Original references}
\makeatletter
\newlabel{algebra-lemma-mod-injective}{{00ME}{source}{}{sourcepointer-1}{}}
\newlabel{algebra-lemma-mod-injective@cref}{{[originalreference][0][]00ME}{[1][0][]source}{}{}{}}
\makeatother
\crefname{originalreference}{original reference}{original references}
\Crefname{originalreference}{Original reference}{Original references}
\makeatletter
\newlabel{coherent-proposition-proper-pushforward-coherent}{{02O5}{source}{}{sourcepointer-2}{}}
\newlabel{coherent-proposition-proper-pushforward-coherent@cref}{{[originalreference][0][]02O5}{[1][0][]source}{}{}{}}
\makeatother
\crefname{originallemma}{original lemma}{original lemmas}
\Crefname{originallemma}{Original lemma}{Original lemmas}
\makeatletter
\newlabel{lemma-h1-fibre-zero-h0-flat}{{0E7K}{source}{}{sourcepointer-3}{}}
\newlabel{lemma-h1-fibre-zero-h0-flat@cref}{{[originallemma][0][]0E7K}{[1][0][]source}{}{}{}}
\makeatother
\crefname{originallemma}{original lemma}{original lemmas}
\Crefname{originallemma}{Original lemma}{Original lemmas}
\makeatletter
\newlabel{lemma-h1-fibre-zero-h0-kappa}{{0E7J}{source}{}{sourcepointer-4}{}}
\newlabel{lemma-h1-fibre-zero-h0-kappa@cref}{{[originallemma][0][]0E7J}{[1][0][]source}{}{}{}}
\makeatother
\crefname{originalreference}{original reference}{original references}
\Crefname{originalreference}{Original reference}{Original references}
\makeatletter
\newlabel{limits-lemma-descend-finite-presentation}{{01ZM}{source}{}{sourcepointer-5}{}}
\newlabel{limits-lemma-descend-finite-presentation@cref}{{[originalreference][0][]01ZM}{[1][0][]source}{}{}{}}
\makeatother
\crefname{originalreference}{original reference}{original references}
\Crefname{originalreference}{Original reference}{Original references}
\makeatletter
\newlabel{limits-lemma-descend-flat-finite-presentation}{{04AI}{source}{}{sourcepointer-6}{}}
\newlabel{limits-lemma-descend-flat-finite-presentation@cref}{{[originalreference][0][]04AI}{[1][0][]source}{}{}{}}
\makeatother
\crefname{originalreference}{original reference}{original references}
\Crefname{originalreference}{Original reference}{Original references}
\makeatletter
\newlabel{limits-lemma-eventually-proper}{{081F}{source}{}{sourcepointer-7}{}}
\newlabel{limits-lemma-eventually-proper@cref}{{[originalreference][0][]081F}{[1][0][]source}{}{}{}}
\makeatother
\crefname{originalreference}{original reference}{original references}
\Crefname{originalreference}{Original reference}{Original references}
\makeatletter
\newlabel{limits-lemma-limit-affine}{{01Z6}{source}{}{sourcepointer-8}{}}
\newlabel{limits-lemma-limit-affine@cref}{{[originalreference][0][]01Z6}{[1][0][]source}{}{}{}}
\makeatother
\title{Vanishing of the first cohomology on the fibre makes the structural map injective}
\author{The Stacks Project Authors}
\date{stacks-project:more-morphisms.tex:local-snapshot-20261011; GNU FDL 1.2 or later}
\begin{document}
\maketitle
\noindent\textit{Editorial note.} Counted span is the author's own proof environment, source lines 21263-21343, immediately following the statement at 21234-21261.\par\medskip
\noindent\textit{Editorial note.} Six cross-chapter pointers (Algebra, Coherent, Limits) are recorded in assembly.external\_references as navigation pointers.\par\medskip
\noindent\textit{Editorial note.} The final reduction to local rings is displayed with the explicit localised map, and the injectivity criterion is applied inside the span.\par\medskip
\noindent\textit{Allowed standard prerequisites (editorial).}\begin{itemize}
\item The flatness statement for H1 vanishing on the fibre. chapter lemma lemma-h1-fibre-zero-h0-flat, cited at proof line 21330. Supplies flatness of the direct image module at the relevant point.
\item The residue field form of the vanishing criterion. chapter lemma lemma-h1-fibre-zero-h0-kappa, cited in the proof. Converts the fibre computation over the residue field into the corresponding statement for the local ring.
\item A module injection criterion. Algebra, lemma algebra-lemma-mod-injective, cited at proof line 21340. Concludes that the map on local rings is injective from flatness plus injectivity of the mod-I reduction.
\item Cohomology of proper pushforward and limit arguments. Coherent, proposition coherent-proposition-proper-pushforward-coherent and the Limits lemmas cited in the proof. Make the direct image coherent and allow the situation to be approximated by finitely presented data.
\end{itemize}\medskip
\ifdefined\hypertarget\hypertarget{sourcepointer-1}{}\fi
\noindent\textit{Source reference.} Original reference 00ME: focused-seeds/geometry/S9-morph-h1-fibre-zero-isom/ref-tags.txt, source lines 4--4. Navigation pointer only: the counted proof cites this statement by \textbackslash{}ref as a standing input; the deduction the counted proof has to make is carried out inside the counted span, and no part of that deduction is delegated to this pointer.\par\medskip
\ifdefined\hypertarget\hypertarget{sourcepointer-2}{}\fi
\noindent\textit{Source reference.} Original reference 02O5: focused-seeds/geometry/S9-morph-h1-fibre-zero-isom/ref-tags.txt, source lines 5--5. Navigation pointer only: the counted proof cites this statement by \textbackslash{}ref as a standing input; the deduction the counted proof has to make is carried out inside the counted span, and no part of that deduction is delegated to this pointer.\par\medskip
\ifdefined\hypertarget\hypertarget{sourcepointer-3}{}\fi
\noindent\textit{Source reference.} Original lemma 0E7K: sources/stacks-more-morphisms/pinned-original-source.tex, source lines 21193--21193. Navigation pointer only: the counted proof cites this statement by \textbackslash{}ref as a standing input; the deduction the counted proof has to make is carried out inside the counted span, and no part of that deduction is delegated to this pointer.\par\medskip
\ifdefined\hypertarget\hypertarget{sourcepointer-4}{}\fi
\noindent\textit{Source reference.} Original lemma 0E7J: sources/stacks-more-morphisms/pinned-original-source.tex, source lines 21059--21059. Navigation pointer only: the counted proof cites this statement by \textbackslash{}ref as a standing input; the deduction the counted proof has to make is carried out inside the counted span, and no part of that deduction is delegated to this pointer.\par\medskip
\ifdefined\hypertarget\hypertarget{sourcepointer-5}{}\fi
\noindent\textit{Source reference.} Original reference 01ZM: focused-seeds/geometry/S9-morph-h1-fibre-zero-isom/ref-tags.txt, source lines 6--6. Navigation pointer only: the counted proof cites this statement by \textbackslash{}ref as a standing input; the deduction the counted proof has to make is carried out inside the counted span, and no part of that deduction is delegated to this pointer.\par\medskip
\ifdefined\hypertarget\hypertarget{sourcepointer-6}{}\fi
\noindent\textit{Source reference.} Original reference 04AI: focused-seeds/geometry/S9-morph-h1-fibre-zero-isom/ref-tags.txt, source lines 7--7. Navigation pointer only: the counted proof cites this statement by \textbackslash{}ref as a standing input; the deduction the counted proof has to make is carried out inside the counted span, and no part of that deduction is delegated to this pointer.\par\medskip
\ifdefined\hypertarget\hypertarget{sourcepointer-7}{}\fi
\noindent\textit{Source reference.} Original reference 081F: focused-seeds/geometry/S9-morph-h1-fibre-zero-isom/ref-tags.txt, source lines 8--8. Navigation pointer only: the counted proof cites this statement by \textbackslash{}ref as a standing input; the deduction the counted proof has to make is carried out inside the counted span, and no part of that deduction is delegated to this pointer.\par\medskip
\ifdefined\hypertarget\hypertarget{sourcepointer-8}{}\fi
\noindent\textit{Source reference.} Original reference 01Z6: focused-seeds/geometry/S9-morph-h1-fibre-zero-isom/ref-tags.txt, source lines 9--9. Navigation pointer only: the counted proof cites this statement by \textbackslash{}ref as a standing input; the deduction the counted proof has to make is carried out inside the counted span, and no part of that deduction is delegated to this pointer.\par\medskip
\begin{lemma}
\label{lemma-h1-fibre-zero-isom}
Consider a commutative diagram
$$
\xymatrix{
X \ar[rr]_f \ar[rd] & & Y \ar[ld] \\
& S
}
$$
of morphisms of schemes. Let $s \in S$ be a point. Assume
\begin{enumerate}
\item $X \to S$ is locally of finite presentation and flat at
points of $X_s$,
\item $Y \to S$ is locally of finite presentation,
\item $f$ is proper,
\item the fibres of $f_s : X_s \to Y_s$ have dimension $\leq 1$
and $R^1f_{s, *}\mathcal{O}_{X_s} = 0$,
\item $\mathcal{O}_{Y_s} \to f_{s, *}\mathcal{O}_{X_s}$ is an isomorphism.
\end{enumerate}
Then there is an open $Y_s \subset V \subset Y$ such that
(a) $V$ is flat over $S$,
(b) $f^{-1}(V)$ is flat over $S$,
(c) $\dim(X_y) \leq 1$ for $y \in V$,
(d) $R^1f_*\mathcal{O}_X|_V = 0$,
(e) $\mathcal{O}_V \to f_*\mathcal{O}_X|_V$
is an isomorphism, and (a) -- (e)
remain true after base change of $f^{-1}(V) \to V$ by any $S' \to S$.
\end{lemma}

\begin{proof}
Let $y \in Y_s$. We may always replace $Y$ by an open neighbourhood of $y$.
Thus we may assume $Y$ and $S$ affine. We may also assume that
$X$ is flat over $S$, $\dim(X_y) \leq 1$ for $y \in Y$,
$R^1f_*\mathcal{O}_X = 0$ universally, and that
$\mathcal{O}_Y \to f_*\mathcal{O}_X$ is surjective, see
Lemma \ref{lemma-h1-fibre-zero-h0-kappa}. (We won't use all of this.)

\medskip\noindent
Assume $S$ and $Y$ affine.
Write $S = \lim S_i$ as a cofiltered of affine Noetherian schemes $S_i$.
By Limits, Lemma \ref{limits-lemma-descend-finite-presentation}
there exists an element $0 \in I$ and a diagram
$$
\xymatrix{
X_0 \ar[rr]_{f_0} \ar[rd] & & Y_0 \ar[ld] \\
& S_0
}
$$
of finite type morphisms of schemes whose base change to $S$
is the diagram of the lemma. After increasing $0$ we may assume $Y_0$
is affine and $X_0 \to S_0$ proper, see
Limits, Lemmas \ref{limits-lemma-eventually-proper} and
\ref{limits-lemma-limit-affine}.
Let $s_0 \in S_0$ be the image of $s$.
As $Y_s$ is affine, we see that
$R^1f_{s, *}\mathcal{O}_{X_s} = 0$ is equivalent
to $H^1(X_s, \mathcal{O}_{X_s}) = 0$.
Since $X_s$ is the base change of $X_{0, s_0}$ by
the faithfully flat map $\kappa(s_0) \to \kappa(s)$
we see that $H^1(X_{0, s_0}, \mathcal{O}_{X_{0, s_0}}) = 0$
and hence $R^1f_{0, *}\mathcal{O}_{X_{0, s_0}} = 0$.
Similarly, as $\mathcal{O}_{Y_s} \to f_{s, *}\mathcal{O}_{X_s}$
is an isomorphism, so is
$\mathcal{O}_{Y_{0, s_0}} \to f_{0, *}\mathcal{O}_{X_{0, s_0}}$.
Since the dimensions of the fibres of $X_s \to Y_s$ are at most $1$,
the same is true for the morphism $X_{0, s_0} \to Y_{0, s_0}$.
Finally, since $X \to S$ is flat, after increasing $0$
we may assume $X_0$ is flat over $S_0$, see
Limits, Lemma \ref{limits-lemma-descend-flat-finite-presentation}.
Thus it suffices to prove the lemma for
$X_0 \to Y_0 \to S_0$ and the point $s_0$.

\medskip\noindent
Combining the reduction arguments above we reduce to the case where
$S$ and $Y$ affine, $S$ Noetherian, the fibres of $f$ have dimension
$\leq 1$, and $R^1f_*\mathcal{O}_X = 0$ universally.
Let $y \in Y_s$ be a point. Claim:
$$
\mathcal{O}_{Y, y} \longrightarrow (f_*\mathcal{O}_X)_y
$$
is an isomorphism. The claim implies the lemma. Namely, since
$f_*\mathcal{O}_X$ is coherent (Cohomology of Schemes, Proposition
\ref{coherent-proposition-proper-pushforward-coherent})
the claim means we can replace $Y$ by an open neighbourhood of $y$
and obtain an isomorphism $\mathcal{O}_Y \to f_*\mathcal{O}_X$.
Then we conclude that $Y$ is flat over $S$
by Lemma \ref{lemma-h1-fibre-zero-h0-flat}.
Finally, the isomorphism $\mathcal{O}_Y \to f_*\mathcal{O}_X$
remains an isomorphism after any base change $S' \to S$
by the final statement of Lemma \ref{lemma-h1-fibre-zero-h0-flat}.

\medskip\noindent
Proof of the claim. We already know that
$\mathcal{O}_{Y, y} \longrightarrow (f_*\mathcal{O}_X)_y$
is surjective (Lemma \ref{lemma-h1-fibre-zero-h0-kappa})
and that $(f_*\mathcal{O}_X)_y$ is
$\mathcal{O}_{S, s}$-flat (Lemma \ref{lemma-h1-fibre-zero-h0-flat})
and that the induced map
$$
\mathcal{O}_{Y_s, y} =  \mathcal{O}_{Y, y}/\mathfrak m_s\mathcal{O}_{Y, y}
\longrightarrow
(f_*\mathcal{O}_X)_y/\mathfrak m_s (f_*\mathcal{O}_X)_y
\to
(f_{s, *}\mathcal{O}_{X_s})_y
$$
is injective by the assumption in the lemma. Then it follows from
Algebra, Lemma \ref{algebra-lemma-mod-injective}
that $\mathcal{O}_{Y, y} \longrightarrow (f_*\mathcal{O}_X)_y$
is injective as desired.
\end{proof}

\end{document}
